\section{한계와 결과 해석 시 주의점}
\label{sec:threats}

\subsection{자료의 한계}
CoC는 사람이 직접 작성한 안전 규칙이 아니다. 주행이 끝난 뒤 기록 영상과 이후 이동 경로를 바탕으로 자동 생성한 설명이다. 따라서 장면의 실제 원인이나 운전 의도를 정확히 설명한다는 보장은 없다. 현재 보관된 변환 자료에는 설명을 만든 모델과 자동 표시 도구의 정확한 버전도 일부 남아 있지 않다. 또한 선택한 27개 장면은 무작위 표본이 아니며 영상과 객체 상태도 제공하지 않으므로, 본 연구는 문장 사이의 관계만 평가했다.

\subsection{평가 범위의 한계}
합성 40쌍은 실제 오류를 따로 수집한 평가 자료가 아니다. 정상 사례에서 한 항목을 의도적으로 바꾸고, 예상한 판정이 나오도록 미리 확인한 시험 사례이다. 따라서 이전 요구를 기억하는 검사의 40/40과 현재 사건만 보는 검사의 20/40은 정해 둔 네 오류 규칙을 얼마나 빠짐없이 실행했는지를 보여 줄 뿐이다. 실제 오류를 맞힌 비율, 실제 오류를 놓치지 않은 비율, 정상 자료를 오류로 잘못 판단한 비율은 아니다. 네 검토자의 첫 판단도 서로 잘 맞지 않았고, 제3자가 모든 합의 과정을 그대로 재현할 기록은 부족하다. 그러므로 검토 후 27/27이 일치했다는 결과도 객관적인 정답으로 볼 수는 없다.

보조 반응 검사는 실제 브레이크나 조향 명령을 사용하지 않고 자차 속도 변화로 반응을 대신 판단했다. 관찰 구간 $W$, 필요한 속도 변화량 $\delta$, 변화 방향의 일관성을 나타내는 비율 $\rho$에 따라 결과가 달라지며, 시험한 27개 조합도 가능한 모든 설정을 포함하지 않는다. 제한 시간 $D=3$초와 STOP/YIELD에서 $v\leq0.5$~m/s이면 이미 요구를 만족했다고 보는 규칙도 별도의 차량 안전 기준에서 가져온 것이 아니다. 특히 후자는 초기 후보를 살펴본 뒤 추가한 탐색 규칙이다. 따라서 $134=125+4+5$는 명시한 설정에서 자료가 어떻게 나뉘었는지를 나타낼 뿐, 차량의 안전률이나 오류율이 아니다.

\subsection{구현과 자료 출처의 한계}
문장에서 정보를 읽는 부분, 입력 자료를 점검하는 부분, 서로 다른 행동 표현을 하나로 맞추는 부분, 결과를 분류하는 부분의 구현 오류가 판정에 영향을 줄 수 있다. \texttt{scene-0001}을 의도적으로 바꾼 시험은 검사기가 미리 넣은 네 오류를 구분하는지만 확인하며, 아직 알려지지 않은 오류까지 찾는 능력은 측정하지 않는다. 현재 자료에는 장면 사이의 기록, 앞뒤 관계, 동일 객체가 이어지는지 보여 주는 정보가 없다. 그래서 기본적으로 각 장면을 서로 독립된 주행 구간으로 처리했다. 장면 경계 시험은 이 처리 방식이 코드에 올바르게 구현됐는지만 확인하며, 실제로 장면 사이에 연속 주행이 없었다는 증거는 아니다.

\subsection{UPPAAL 모델과 일반화의 한계}
UPPAAL 모델은 미리 정한 사건 순서를 그대로 재생하고 Python 검사기와 같은 상태 변경 규칙을 사용한다. 7개 시험 사례에서 결과가 모두 같았다는 사실은 두 구현이 해당 사례에서 일치했음을 보여 주지만, 규칙 자체가 현실을 올바르게 표현한다는 독립적인 증명은 아니다. 상대 객체의 상태, 실제 시간 제한, 브레이크ㆍ조향 명령, 차량의 물리적 움직임과 제어 결과가 다시 차량에 영향을 주는 과정도 포함하지 않았다. 따라서 다른 자료ㆍ언어ㆍCoC 생성 방식으로의 일반화, 실제 차량 안전, 제어 명령 생성 또는 현장 적용 가능성을 입증하지 않는다.
